\section{Обсуждение}

\subsection{Сравнение подходов к обеспечению идемпотентности}

Результаты исследования показывают, что каждый тип контроллера Spring требует различного подхода к обеспечению идемпотентности (таблица~\ref{tab:comparison}).

\begin{table}[H]
\centering
\caption{Сравнение механизмов идемпотентности по типам контроллеров}
\label{tab:comparison}
\small
\begin{tabularx}{\textwidth}{@{}l X X X@{}}
\toprule
\textbf{Критерий} & \textbf{@RestController} & \textbf{Functional} & \textbf{Data REST} \\
\midrule
Механизм POST & Idempotency-Key + фильтр & Idempotency-Key + фильтр & Unique constraints \\
Механизм PUT & Естественная идемпотентность & Естественная идемпотентность & @Version + ETag \\
Concurrent защита & На уровне фильтра & На уровне фильтра & Optimistic locking \\
Сложность реализации & Средняя & Средняя & Низкая (автоматически) \\
Гибкость & Высокая & Высокая & Ограниченная \\
\bottomrule
\end{tabularx}
\end{table}

\subsection{Анализ результатов}

\textbf{Аннотационный \mbox{@RestController}} предоставляет наибольшую гибкость в реализации идемпотентности, но требует ручной реализации фильтра и хранилища ключей. Этот подход наиболее распространён в production-системах благодаря явности и читаемости кода. Разработчик полностью контролирует обработку запроса, может добавить произвольную логику валидации, трансформации и обработки ошибок. Недостаток~--- необходимость писать и поддерживать больше кода.

\textbf{Функциональные эндпоинты} (WebMvc.fn) используют тот же механизм Idempotency-Key через общий фильтр, что демонстрирует важное архитектурное преимущество --- фильтры работают на уровне Servlet API, независимо от типа контроллера. Это означает, что один и тот же \texttt{IdempotencyFilter} обеспечивает идемпотентность как для \mbox{\texttt{@RestController}}, так и для \texttt{RouterFunction}. Функциональный стиль обеспечивает лучшую композиционность: маршруты можно комбинировать, вкладывать и переиспользовать программно, в отличие от аннотационного подхода, где маршруты привязаны к методам классов.

\textbf{Spring Data REST} автоматически обеспечивает ряд механизмов идемпотентности: \texttt{@Version} генерирует ETag, а unique constraints предотвращают дубликаты. Это минимизирует объём ручного кода --- разработчику достаточно определить JPA-сущность и репозиторий. Однако этот подход ограничивает возможности кастомизации: Idempotency-Key не поддерживается автоматически, а бизнес-логика добавляется через EventHandler, что менее гибко, чем прямое управление в Service-слое.

\subsection{Значимость Testcontainers в тестировании идемпотентности}

Использование Testcontainers вместо in-memory баз данных (H2) имеет критическое значение для тестирования идемпотентности. Ряд механизмов зависит от особенностей конкретной СУБД:

\begin{itemize}
    \item \textbf{UNIQUE constraints} --- поведение при нарушении уникальности различается между H2 и PostgreSQL (разные коды ошибок, разная обработка транзакций).
    \item \textbf{Flyway-миграции} --- синтаксис SQL может отличаться (например, \texttt{gen\_random\_uuid()} специфичен для PostgreSQL).
    \item \textbf{Конкурентный доступ} --- PostgreSQL поддерживает уровни изоляции и блокировки, отличные от H2.
\end{itemize}

Singleton-паттерн Testcontainers (один контейнер на все тесты) обеспечивает баланс между точностью и скоростью: тесты работают с реальной СУБД, но контейнер запускается только один раз за сессию.

\subsection{Рекомендации для production}

На основе проведённого исследования сформулированы следующие рекомендации:

\begin{enumerate}
    \item \textbf{Всегда требовать Idempotency-Key для POST-запросов}, создающих ресурсы или выполняющих побочные эффекты. Это предотвращает дублирование при retry и является отраслевым стандартом~\cite{stripe-idempotency}.
    \item \textbf{Использовать optimistic locking} (\texttt{@Version}) для защиты от конкурентных обновлений. Это особенно важно для систем с высокой нагрузкой.
    \item \textbf{Устанавливать TTL} для ключей идемпотентности. Рекомендуемое значение --- 24 часа~\cite{ietf-idempotency-key}. Слишком короткий TTL не защитит от поздних retry, слишком длинный --- приведёт к росту таблицы ключей.
    \item \textbf{Хранить ключи} в той же СУБД, что и бизнес-данные, для обеспечения транзакционной консистентности. Альтернатива (Redis) обеспечивает лучшую производительность, но вводит риск рассогласования.
    \item \textbf{Тестировать с реальной СУБД} через Testcontainers~\cite{testcontainers}, а не с in-memory заменителями, для достоверности результатов.
    \item \textbf{Комбинировать механизмы}: Idempotency-Key для POST, @Version для PUT, UNIQUE constraints как последний рубеж защиты от дубликатов.
\end{enumerate}

\section{Заключение}

В данной работе проведено исследование и демонстрация тестирования идемпотентности REST API на базе Spring Boot. Разработано приложение с тремя кардинально различными типами контроллеров, для каждого из которых реализованы и протестированы механизмы обеспечения идемпотентности.

Основные результаты:
\begin{itemize}
    \item Формализована математическая модель идемпотентности HTTP-запросов с доказательством корректности паттерна Idempotency-Key.
    \item Реализован и протестирован паттерн Idempotency-Key с хранением ключей в PostgreSQL, применимый к любому типу контроллера через Servlet-фильтр.
    \item Продемонстрированы три кардинально различных подхода к организации REST API в Spring: аннотационный (@RestController), функциональный (RouterFunction) и автогенерируемый (Spring Data REST), каждый с собственным механизмом обеспечения идемпотентности.
    \item Разработано 22 интеграционных теста с использованием Testcontainers и PostgreSQL, покрывающих все аспекты идемпотентности для трёх типов контроллеров.
    \item Продемонстрировано критическое значение механизмов идемпотентности: при их отключении (профиль non-idempotent) система допускает дублирование данных и тройное списание при retry.
\end{itemize}
